\section{Polynomials over Finite Fields}

Throughout this section, $q$ is a prime power and $\F_q$ is the field with $q$ elements.

\subsection{Sums over $\F_q$}

Summing a power of $x$ over all of $\F_q$ gives zero, unless the power is large.

% [CORRECTED] was: "$k \leq q - 1$" --- at $k = q - 1$ the sum is $q - 1 = -1$; and the proof below was: "$x^{k} - 1$ has at most $k < q$ roots over $\F_q$. So there is certainly some $a \in \F_q$ that is" --- that $a$ could be $0$, and $x^0 - 1 = 0$
\begin{boxtheorem}
    For all $k < q - 1$,
    \begin{align*}
        \sum_{x \in \F_q} x^k = 0
    \end{align*}
\end{boxtheorem}
\begin{proof}
    For $k = 0$ the sum is $q = 0$. Otherwise, note that $x^{k} - 1$ has at most $k < q - 1$ roots over $\F_q$. So there is certainly some $a \in \F_q \setminus \set{0}$ that is \textit{not} a root of this polynomial. So $a^k \neq 1$. Since $x \mapsto ax$ is a bijection,
    \begin{align*}
        \sum_{x \in \F_q} x^k = \sum_{x \in \F_q} \parenth{ax}^k = a^k \sum_{x \in \F_q} x^k
    \end{align*}
    As $a^k \ne 1$, the only way this equality can hold is if $\sum_{x \in \F_q} x^k = 0$.
\end{proof}

Summing over $\F_q^n$ instead, the same is true of any polynomial of small enough degree.

\begin{boxtheorem} \label{Ch8:Thm:Sum_Over_Fq}
    Let $p \in \F_q\!\brac{x_1, \ldots, x_n}$. Assume $\deg(p) < n(q - 1)$. Then
    \begin{align*}
        \sum_{\bx \in \F_q^n} p\of{\bx} = 0
    \end{align*}
\end{boxtheorem}
\begin{proof}
    Enough to do the case where $p(\bx)$ is a monomial. \sorry
    % Hint: just fix j_1, \ldots, j_n whose sum is less than n(q - 1)
\end{proof}

\subsection{Fermat's Little Theorem}

We'll now prove Fermat's Little Theorem combinatorially.

\begin{boxtheorem}[Fermat's Little Theorem]
    For $p$ a prime, for all $a \in \Z$, $a^p \equiv a \pmod{p}$. Equivalently, for $p$ a prime, for all $a \in \Z$ such that $p \nmid a$, $a^{p - 1} \equiv 1 \pmod{p}$.
\end{boxtheorem}
\begin{proof}
    % [CORRECTED] was: "We know that $x \mapsto ax$ is a bijection", and the display ended "$= a^{p - 1}$" --- the bijection needs $p \nmid a$, and the product is $-1$, not $1$, so it is cancelled rather than dropped
    For $p \nmid a$, we know that $x \mapsto ax$ is a bijection of $\F_p \setminus \set{0}$. So
    \begin{align*}
        \prod_{x \in \F_p \setminus \set{0}} x
        = \prod_{x \in \F_p \setminus \set{0}} (ax)
        = a^{p - 1} \prod_{x \in \F_p \setminus \set{0}} x
    \end{align*}
    in $\F_p$, and cancelling the product, which is nonzero, gives $a^{p - 1} = 1$.
\end{proof}

Of course the corresponding result holds in $\F_q$ as well.

\begin{boxtheorem}
    For all $a \in \F_q \setminus \set{0}$, $a^{q - 1} = 1$.
\end{boxtheorem}
\begin{proof}[``Proof'']
    $\F_q$ is, by definition, the splitting field of $x^{q} - x$ over $\F_p$, where $q$ is a power of $p$. So when $a \in \F_q$ is nonzero,
    \begin{align*}
        a^q - a = 0 \iff a\parenth{a^{q - 1} - 1} = 0 \iff a^{q - 1} - 1 = 0
    \end{align*}
\end{proof}

\subsection{The Chevalley--Warning Theorem}

Let's now study systems of polynomials over $\F_q$.

\begin{boxtheorem}[Chevalley-Warning]
    Fix $p_1, \ldots, p_k \in \F_q\!\brac{x_1, \ldots, x_n}$ with $\sum_{i=1}^{k} \pdeg{p_i} < n$. Let $p$ be the characteristic of $\F_q$. Then
    \begin{align*}
        p \mid \#\!\setst{\bx \in \F_q^n}{p_1\of{\bx} = p_2\of{\bx} = \cdots = p_k\of{\bx} = 0}
    \end{align*}
\end{boxtheorem}
\begin{proof}
    Define the polynomial
    \begin{align*}
        Q(\bx) = \prod_{i=1}^{k} \parenth{1 - p_i(\bx)^{q - 1}}
    \end{align*}
    What's interesting about the way we've defined $Q$ is that
    \begin{itemize}
        \item $Q(\bx) = 0$ iff $p_i(\bx) \ne 0$ for \textit{some} $i$
        \item $Q(\bx) = 1$ iff $p_i(\bx) = 0$ for \textit{every} $i$
    \end{itemize}
    Consider the quantity
    \begin{align*}
        \F_q \ni z = \sum_{\bx \in \F_q^n} Q(\bx)
    \end{align*}
    Notice that $\pdeg{Q} < n(q - 1)$, so by \Cref{Ch8:Thm:Sum_Over_Fq}, 
    $z = 0$. We can write
    \begin{align*}
        z = \underbrace{1 + \cdots + 1}_{N \text{ times}}
    \end{align*}
    where $N = \#\!\setst{\bx \in \F_q^n}{p_1\of{\bx} = p_2\of{\bx} = \cdots = p_k\of{\bx} = 0}$. As $N = 0$ in $\F_q$, we must have $p \mid N$.
\end{proof}

\subsection{The \protect\Erdos--Ginzburg--Ziv Theorem}

Chevalley--Warning has a striking consequence for the integers.

\begin{boxtheorem}[\Erdos-Ginzburg-Ziv]
    Out of any $2n - 1$ integers, there are $n$ whose average is an integer. More precisely, for all $a_1, \ldots, a_{2n - 1} \in \Z$, $\exists I \ssq \set{1, \ldots, 2n - 1}$ such that $\abs{I} = n$ and $\sum_{i \in I} a_i \equiv 0 \pmod{n}$.
\end{boxtheorem}

\begin{proof}
    % [CORRECTED] was: the display "$p(\bx) = \sum_{\bx \in \F_p^{2n - 1}} \indicator_I(x)$" --- its right-hand side sums out $\bx$, and $p$ is the prime here; this is the identity the argument below uses, by Fermat's Little Theorem
    Consider the case $n = p$, with $p$ a prime. Notice that if $\bx = \parenth{x_1, \ldots, x_{2p - 1}} \in \F_p^{2p - 1}$,
    \begin{align*}
        \sum_{i=1}^{2p - 1} x_i^{p - 1} = \sum_{i=1}^{2p - 1} \indicator_{x_i \neq 0}
    \end{align*}
    % [FILLED] the prime case, from "Fix $a_1, \ldots$" to the end of this paragraph, was supplied; it is the standard deduction from Chevalley--Warning
    Fix $a_1, \ldots, a_{2p - 1} \in \Z$, and consider the polynomials
    \begin{align*}
        p_1\of{\bx} = \sum_{i=1}^{2p - 1} a_i x_i^{p - 1} \qquad \text{and} \qquad p_2\of{\bx} = \sum_{i=1}^{2p - 1} x_i^{p - 1}
    \end{align*}
    in $\F_p\!\brac{x_1, \ldots, x_{2p - 1}}$, reading the $a_i$ mod $p$. Their degrees add up to at most $2p - 2 < 2p - 1$, so by Chevalley--Warning, the number of common zeros of $p_1$ and $p_2$ is a multiple of $p$. One of them is $\bx = 0$, so there is another, $\bx \neq 0$. Let $I = \setst{i}{x_i \neq 0}$, which is nonempty. By Fermat's Little Theorem, $x_i^{p - 1}$ is $1$ if $i \in I$ and $0$ otherwise, so in $\F_p$,
    \begin{align*}
        0 = p_2\of{\bx} = \abs{I} \qquad \text{and} \qquad 0 = p_1\of{\bx} = \sum_{i \in I} a_i
    \end{align*}
    So $p$ divides $\abs{I}$, and as $1 \leq \abs{I} \leq 2p - 1$, $\abs{I} = p$. And $\sum_{i \in I} a_i \equiv 0 \pmod{p}$.

    This takes care of the prime case.

    % [CORRECTED] was: "for $n = a$ and $n = b$, with $a$ and $b$ coprime", and below, "Using the result for $a$, I can find $a$ of the sets $B_{i_1}, B_{i_2}, \ldots, B_{i_a}$ such that $\sum_{j=1}^{a} y_{ij} \equiv 0 \pmod{a}$. So $\sum_{y_{ij}} \equiv 0 \pmod{a}$." (which phase 2 of the post-lecture pass finished using coprimality) --- from the primes, coprime products reach only squarefree $n$ (never $4$); applying the result for $a$ to the $\frac{y_i}{b}$ needs no coprimality
    Now, assume the result is proven for $n = a$ and $n = b$. We prove the result for $n = ab$.

    Fix $c_1, c_2, \ldots, c_{2ab - 1} \in \Z$. I repeatedly find and remove $b$ terms summing to $0$ mod $b$. I can do this $2a - 1$ times, to get disjoint sets $B_1, B_2, \ldots, B_{2a - 1}$, each summing as $\sum_{c \in B_i} c = y_i \equiv 0 \pmod{b}$.

    Using the result for $a$ on the integers $\frac{y_1}{b}, \ldots, \frac{y_{2a - 1}}{b}$, I can find $a$ of the sets $B_{i_1}, B_{i_2}, \ldots, B_{i_a}$ such that $\sum_{j=1}^{a} \frac{y_{i_j}}{b} \equiv 0 \pmod{a}$, so that $\sum_{j=1}^{a} y_{i_j} \equiv 0 \pmod{ab}$. But this is the sum of the $ab$ numbers in $B_{i_1} \cup \cdots \cup B_{i_a}$, which is what we wanted.
\end{proof}
